% Capitulo: Scrum com Kanban: o modelo que usei
\chapter{Scrum com Kanban: o modelo que usei}
\label{cap:scrumban}

\rascunho{
  \textbf{Ideia do capítulo.} O modelo híbrido aplicado nos meus projetos, descrito de forma que outro time consiga copiar e adaptar \parencite{ladas2009}.
  \begin{itemize}
    \item Cadências: o que ficou do Scrum (planejamento, revisão, retrospectiva) e com que frequência.
    \item O quadro: colunas, limite de trabalho em andamento (WIP), classes de serviço.
    \item Refinamento de demandas de dados: o que uma história precisa ter para entrar.
      Com IA, os testes de aceite podem ser escritos já no refinamento, junto com o
      stakeholder, e a história entra no quadro com a regra de negócio em forma de teste
      (ver \veja{sec:tdd}).
    \item Definição de pronto para dados e para dashboards: testes escritos antes e
      passando como critério obrigatório, não opcional.
    \item A frequência de atualização é definida no refinamento, a partir do ritmo da
      decisão, e revisada com as métricas de uso (ver o capítulo~\ref{cap:frequencia}).
    \item Métricas do fluxo: tempo de ciclo, vazão.
  \end{itemize}
}

\section{Cadências}
\label{sec:cadencias}

\pendente{escrever}

\section{O quadro e o limite de WIP}
\label{sec:quadro-wip}

\pendente{escrever}

\section{Refinamento}
\label{sec:refinamento}

\pendente{escrever}

\section{Definição de pronto}
\label{sec:definicao-pronto}

\pendente{escrever}

\section{Métricas do fluxo}
\label{sec:metricas-fluxo}

\pendente{escrever}
